%% ma: L12-B14-0038
%% lop: 12
%% chuong: 5
%% bai: 14
%% dang: 12.14.19
%% loai: TN4
%% muc: TH
%% dap_an: "A"
%% nguon: "CĐ2.5.pdf giáo viên gửi 27/10/2026 (M07, Lê Minh Khôi – Nguyễn Anh Khôi, Chương II Chủ đề 5 Tích có hướng của hai vectơ), Câu 22"
%% kien_thuc: "Bài 14 (tích có hướng, phương trình mặt phẳng, khoảng cách từ điểm đến mặt phẳng), lớp 11 (thể tích khối chóp, khối hộp)"
%% trang_thai: chinh-thuc
%% ghi_chu: "Giáo viên đã duyệt 27/10/2026"
\begin{ex}
Trong không gian $Oxyz$, cho bốn điểm $A(0;0;2)$, $B(3;0;5)$, $C(1;1;0)$, $D(4;1;2)$. Độ dài đường cao của tứ diện $ABCD$ hạ từ đỉnh $D$ xuống mặt phẳng $(ABC)$ bằng
\choice{\True $\dfrac1{\sqrt{11}}$}{$1$}{$11$}{$\sqrt{11}$}
\loigiai{$\left[\overrightarrow{AB},\overrightarrow{AC}\right]=(-3;9;3)$ nên $(ABC)$: $x-3y-z+2=0$.

$d\left(D;(ABC)\right)=\dfrac{|4-3-2+2|}{\sqrt{11}}=\dfrac1{\sqrt{11}}$.}
\end{ex}
